 \documentclass[opr,utf8]{sinol}
 
 \usepackage{tikz}
 
  \title{Obliczenia w chmurze}
  \id{obl}
  \signature{kpok???}
  \author{Karol Pokorski} %% autor opracowania
  \history{2.08.2018}{KPok, przygotowanie opracowania autorskiego}{1.00}

\begin{document}
  \begin{tasktext}%

  \section{Rozwiązanie wzorcowe}
    Wrzucamy oferty procesorów i~zadania do~jednego wora.
    Sortujemy to~względem częstotliwości malejąco.

    Przetwarzamy teraz zdarzenia dwóch typów:
    \begin{itemize}
      \item oferta procesora, który wystarczy częstotliwością taktowania
      na~wykonanie wszystkich zadań, które przeanalizujemy za~chwilę,
      \item zlecenie obliczenia.
    \end{itemize}
    Tym samym (dzięki sortowaniu) wyeliminowaliśmy problem zajmowania się
    częstotliwością taktowania.

    Teraz jak przetwarzać zdarzenia.
    Utrzymujemy dynamika, którego stanem jest liczba już kupionych, ale jeszcze
    nieużytych rdzeni procesorów.

    Gdy dostajemy ofertę procesora aktualizujemy tablicę dynamika liniowo względem
    liczby stanów, podobnie jak w~problemie plecakowym.
    Analogicznie (tylko od~drugiej strony) postępujemy z~zadaniami do~wykonania.

    Złożoność to: łączna liczba rdzeni wszystkich procesorów razy liczba zdarzeń.
    Z~uwagi na~naprawdę małą stałą, możnaby nawet jeszcze podbić te limity
    co~są w~zadaniu (powiedzmy $n, m \le 5\,000$ na~moim kompie działa sekundę).
    Nie wiem czy chcemy, zawodnicy mogą nie zgadnąć, że~to będzie dostatecznie
    szybkie.

    Wzorcówka jest już zaimplementowana w~\solution{.cpp}.

  \section{Potencjalne inne rozwiązania}
    \begin{itemize}
      \item zgadywanie podzbioru procesorów do~kupienia i~ustalanie potem
      zadań, które dadzą się zrobić (nie wiem czy się da),
      \item zgadywanie podzbioru zadań do~wykonania i~ustalanie potem
      procesorów, które należy kupić (nie wiem czy się da),
      \item randomowe shufflowanie kolejności zadań i procesorów, kupowanie
      prefiksu listy procesorów, aby wykonać prefiks listy zadań (wybór optymalnego
      prefiksu w optymalnym random shufflowaniu)
      \item jakieś zachłany oparte o jakieś dziwne współczynniki,
      \item rozwiązanie bez obserwacji o~sortowaniu po~częstotliwości (działa, gdy $f_i = F_i = 1$),
      \item matching ważony między procesorami a~zadaniami (działa, gdy $c_i = C_i = 1$),
    \end{itemize}

  \section{Do rozważenia}
    Być może chcemy wymagać od~zawodników wypisania pełnej odpowiedzi (tj.
    listy kupionych procesorów i~wykonanych zadań)... jeśli Opracowujący idzie
    w~ten wariant, niech upewni się, że zarówno zmodyfikowana wzorcówka jak
    i~weryfikator outputa jest poprawny. W~razie wątpliwości, lepiej się nie męczyć.

  \end{tasktext}
\end{document}
